\ifdefined\tsenvDocsCombinedAppendixMode
\def\tsenvAppendixPromptExampleCombinationsStart{%
  \docTitle{Appendix: Prompt Example Combinations}%
  \phantomsection\label{docstart:appendix_prompt_example_combinations}%
  \section*{Prompt Example Combinations}%
}
\def\tsenvAppendixPromptExampleCombinationsEnd{}
\else\ifdefined\tsenvCombinedAppendixMode
\def\tsenvAppendixPromptExampleCombinationsStart{%
  \section{Prompt Example Combinations}%
  \label{app:prompt-example-combinations}%
}
\def\tsenvAppendixPromptExampleCombinationsEnd{}
\else
\documentclass[11pt]{article}
\usepackage{manual_docs_style}

\def\tsenvAppendixPromptExampleCombinationsStart{%
  \begin{document}%
  \docTitle{Appendix: Prompt Example Combinations}%
  \phantomsection\label{docstart:appendix_prompt_example_combinations}%
  \section*{Prompt Example Combinations}%
}
\def\tsenvAppendixPromptExampleCombinationsEnd{\end{document}}
\fi\fi
\tsenvAppendixPromptExampleCombinationsStart

The high- and low-description prompts are derived from \texttt{BallDrop}; the none-description prompts intentionally hide the domain semantics.
Each subsection shows a different prompt case, named with the convention \code{<description\_level>\_<task\_type>\_<data\_setting>}, where \code{data\_setting} is either \code{0} or \code{n}.

\subsection*{\code{high_desc_direct_0}}
\begin{DocCode}
The time series in the test_samples/ folder were generated by simulating a 1D vertical point-mass ball model under gravity.
While the ball is above the ground, its motion is governed by gravity and quadratic air drag acting opposite the direction of travel, and the position evolves according to the current velocity.
Ground interaction is modeled with a restitution-based hard-stop law.
When the ball reaches the ground with impact speed at or above 0.5 m/s, the impact is treated as instantaneous and the post-impact speed is set by the coefficient of restitution e in [0,1].
The rebound points upward and its magnitude equals e times the pre-impact speed.
When the impact speed is below the threshold, the ball does not rebound and instead enters static contact at the ground.

Observed Signals:
col1: ball height
col2: ball velocity
col3: contact impulse (N*s)
col4: time

For each simulation, either no parameter changes occur, or exactly one parameter among the allowed labels changes during the observed simulation interval.
If a parameter changes, it undergoes a single instantaneous step change at an unknown time during the observed interval.

Allowed labels:
["coefficient of restitution", "mass", "drag coefficient", "gravity acceleration", "no parameter change"]

Use "no parameter change" if there is no evidence in the data of a parameter change during the observed interval.

Task:
Create a file named results.json in the current working directory.
For each file in test_samples/, return a ranked list of labels.
The first label is your final top-1 prediction and should be the single label you think is most likely correct.
You may include additional labels only when the evidence is genuinely ambiguous.
Additional labels are treated as lower-confidence alternatives.
The output must be valid JSON with exactly this structure:
{

  "<filename_1>": ["<top_1_label>"],
  "<filename_2>": ["<top_1_label>", "<optional_lower_confidence_label>"]

}

Requirements:
- Include one entry for every Parquet file in test_samples/.
- Every returned label must exactly match one of the allowed labels.
- The order of labels matters: the first label is the top-1 prediction.
- Do not include duplicate labels for a sample.
- Do not return all labels unless the evidence is genuinely ambiguous across all labels.

Evaluation:
The primary evaluation metric is top-1 accuracy: the first label in each returned list is compared with the hidden correct label.
A secondary shortlist score may also be reported. For a returned list of length m, the sample receives score 1/m if the hidden correct label appears anywhere in the list, and 0 otherwise. Therefore, unnecessary extra labels reduce the secondary score.

Additional requirements:
- If you create intermediate files, images, scripts, or notes while solving the task, create them in the current working directory.
- Internet access is disabled.
\end{DocCode}


\subsection*{\code{high_desc_code_0}}
\begin{DocCode}
Context:
The time series in the test_samples/ folder were generated by simulating a 1D vertical point-mass ball model under gravity.
While the ball is above the ground, its motion is governed by gravity and quadratic air drag acting opposite the direction of travel, and the position evolves according to the current velocity.
Ground interaction is modeled with a restitution-based hard-stop law.
When the ball reaches the ground with impact speed at or above 0.5 m/s, the impact is treated as instantaneous and the post-impact speed is set by the coefficient of restitution e in [0,1].
The rebound points upward and its magnitude equals e times the pre-impact speed.
When the impact speed is below the threshold, the ball does not rebound and instead enters static contact at the ground.

Observed Signals:
col1: ball height
col2: ball velocity
col3: contact impulse (N*s)
col4: time

For each simulation, either no parameter changes occur, or exactly one parameter among the allowed labels changes during the observed simulation interval.
If a parameter changes, it undergoes a single instantaneous step change at an unknown time during the observed interval.

Allowed labels:
["coefficient of restitution", "mass", "drag coefficient", "gravity acceleration", "no parameter change"]

Use "no parameter change" if there is no evidence in the data of a parameter change during the observed interval.

Task:
Create a Python script named rule.py in the current working directory.
The script must define exactly this function:
def predict(df) -> list[str]:
The input df is a pandas DataFrame containing one sample with columns col1, col2, col3, and col4.
The function will be called on samples inside test_samples/ and on additional held-out samples with the same schema and label set.
For each dataframe, predict(df) must return a ranked list of labels. The first label is the final top-1 prediction and should be the single label most likely to be correct. Additional labels are optional lower-confidence alternatives and should only be included when the evidence is genuinely ambiguous.

Requirements for predict(df):
- Return a Python list of strings.
- Every returned label must exactly match one of the allowed labels.
- The first returned label is the top-1 prediction.
- Do not include duplicate labels.
- Do not return all labels unless the evidence is genuinely ambiguous across all labels.
- You may inspect any file while developing rule.py, but the final submitted rule.py must not read, open, import, or depend on any files at prediction time, including test_samples/
- The final rule.py must be able to run on a dataframe alone.

Evaluation:
The primary evaluation metric is top-1 accuracy: the first label returned by predict(df) is compared with the hidden correct label.
A secondary shortlist score may also be reported. For a returned list of length m, the sample receives score 1/m if the hidden correct label appears anywhere in the list, and 0 otherwise. Therefore, unnecessary extra labels reduce the secondary score.

Additional requirements:
- If you create intermediate files, images, scripts, or notes while solving the task, create them in the current working directory.
- Internet access is disabled.
\end{DocCode}


\subsection*{\code{high_desc_direct_n}}
\begin{DocCode}
The time series in the test_samples/ and train_samples/ folders were generated by simulating a 1D vertical point-mass ball model under gravity.
While the ball is above the ground, its motion is governed by gravity and quadratic air drag acting opposite the direction of travel, and the position evolves according to the current velocity.
Ground interaction is modeled with a restitution-based hard-stop law.
When the ball reaches the ground with impact speed at or above 0.5 m/s, the impact is treated as instantaneous and the post-impact speed is set by the coefficient of restitution e in [0,1].
The rebound points upward and its magnitude equals e times the pre-impact speed.
When the impact speed is below the threshold, the ball does not rebound and instead enters static contact at the ground.

Observed Signals:
col1: ball height
col2: ball velocity
col3: contact impulse (N*s)
col4: time

For each simulation, either no parameter changes occur, or exactly one parameter among the allowed labels changes during the observed simulation interval.
If a parameter changes, it undergoes a single instantaneous step change at an unknown time during the observed interval.

Allowed labels:
["coefficient of restitution", "mass", "drag coefficient", "gravity acceleration", "no parameter change"]

Use "no parameter change" if there is no evidence in the data of a parameter change during the observed interval.

Task:
Create a file named results.json in the current working directory.
For each file in test_samples/, return a ranked list of labels.
The first label is your final top-1 prediction and should be the single label you think is most likely correct.
You may include additional labels only when the evidence is genuinely ambiguous.
Additional labels are treated as lower-confidence alternatives.
The output must be valid JSON with exactly this structure:
{

  "<filename_1>": ["<top_1_label>"],
  "<filename_2>": ["<top_1_label>", "<optional_lower_confidence_label>"]

}

To help with this task, you can use the labeled `train_samples/` directory. The corresponding labels are available in `train_labels.json` file.

Requirements:
- Include one entry for every Parquet file in test_samples/.
- Every returned label must exactly match one of the allowed labels.
- The order of labels matters: the first label is the top-1 prediction.
- Do not include duplicate labels for a sample.
- Do not return all labels unless the evidence is genuinely ambiguous across all labels.

Evaluation:
The primary evaluation metric is top-1 accuracy: the first label in each returned list is compared with the hidden correct label.
A secondary shortlist score may also be reported. For a returned list of length m, the sample receives score 1/m if the hidden correct label appears anywhere in the list, and 0 otherwise. Therefore, unnecessary extra labels reduce the secondary score.

Additional requirements:
- If you create intermediate files, images, scripts, or notes while solving the task, create them in the current working directory.
- Internet access is disabled.
\end{DocCode}

\subsection*{\code{high_desc_code_n}}
\begin{DocCode}
Context:
The time series in the test_samples/ and train_samples/ folders were generated by simulating a 1D vertical point-mass ball model under gravity.
While the ball is above the ground, its motion is governed by gravity and quadratic air drag acting opposite the direction of travel, and the position evolves according to the current velocity.
Ground interaction is modeled with a restitution-based hard-stop law.
When the ball reaches the ground with impact speed at or above 0.5 m/s, the impact is treated as instantaneous and the post-impact speed is set by the coefficient of restitution e in [0,1].
The rebound points upward and its magnitude equals e times the pre-impact speed.
When the impact speed is below the threshold, the ball does not rebound and instead enters static contact at the ground.

Observed Signals:
col1: ball height
col2: ball velocity
col3: contact impulse (N*s)
col4: time

For each simulation, either no parameter changes occur, or exactly one parameter among the allowed labels changes during the observed simulation interval.
If a parameter changes, it undergoes a single instantaneous step change at an unknown time during the observed interval.

Allowed labels:
["coefficient of restitution", "mass", "drag coefficient", "gravity acceleration", "no parameter change"]

Use "no parameter change" if there is no evidence in the data of a parameter change during the observed interval.

Task:
Create a Python script named rule.py in the current working directory.
The script must define exactly this function:
def predict(df) -> list[str]:
The input df is a pandas DataFrame containing one sample with columns col1, col2, col3, and col4.
The function will be called on samples inside test_samples/ and on additional held-out samples with the same schema and label set.
For each dataframe, predict(df) must return a ranked list of labels. The first label is the final top-1 prediction and should be the single label most likely to be correct. Additional labels are optional lower-confidence alternatives and should only be included when the evidence is genuinely ambiguous.

To help with this task, you can use the labeled train_samples/ directory while developing rule.py. The corresponding labels are available in train_labels.json.

Requirements for predict(df):
- Return a Python list of strings.
- Every returned label must exactly match one of the allowed labels.
- The first returned label is the top-1 prediction.
- Do not include duplicate labels.
- Do not return all labels unless the evidence is genuinely ambiguous across all labels.
- You may inspect train_samples/ and train_labels.json while developing rule.py, but the final submitted rule.py must not read, open, import, or depend on any files at prediction time, including train_samples/, test_samples/ or train_labels.json
- The final rule.py must be able to run on a dataframe alone.

Evaluation:
The primary evaluation metric is top-1 accuracy: the first label returned by predict(df) is compared with the hidden correct label.
A secondary shortlist score may also be reported. For a returned list of length m, the sample receives score 1/m if the hidden correct label appears anywhere in the list, and 0 otherwise. Therefore, unnecessary extra labels reduce the secondary score.

Additional requirements:
- If you create intermediate files, images, scripts, or notes while solving the task, create them in the current working directory.
- Internet access is disabled.
\end{DocCode}

\subsection*{\code{low_desc_direct_n}}
\begin{DocCode}
The time series in the test_samples/ and train_samples/ folders were generated by simulating a ball bouncing on a surface.

Observed Signals:
col1: ball position (height)
col2: ball velocity
col3: contact impulse (N*s)
col4: time

For each simulation, either no parameter changes occur, or exactly one parameter among the allowed labels changes during the observed simulation interval.
If a parameter changes, it undergoes a single instantaneous step change at an unknown time during the observed interval.

Allowed labels:
["coefficient of restitution", "mass", "drag coefficient", "gravity acceleration", "no parameter change"]

Use "no parameter change" if there is no evidence in the data of a parameter change during the observed interval.

Task:
Create a file named results.json in the current working directory.
For each file in test_samples/, return a ranked list of labels.
The first label is your final top-1 prediction and should be the single label you think is most likely correct.
You may include additional labels only when the evidence is genuinely ambiguous.
Additional labels are treated as lower-confidence alternatives.
The output must be valid JSON with exactly this structure:
{

  "<filename_1>": ["<top_1_label>"],
  "<filename_2>": ["<top_1_label>", "<optional_lower_confidence_label>"]

}

To help with this task, you can use the labeled `train_samples/` directory. The corresponding labels are available in `train_labels.json` file.

Requirements:
- Include one entry for every Parquet file in test_samples/.
- Every returned label must exactly match one of the allowed labels.
- The order of labels matters: the first label is the top-1 prediction.
- Do not include duplicate labels for a sample.
- Do not return all labels unless the evidence is genuinely ambiguous across all labels.

Evaluation:
The primary evaluation metric is top-1 accuracy: the first label in each returned list is compared with the hidden correct label.
A secondary shortlist score may also be reported. For a returned list of length m, the sample receives score 1/m if the hidden correct label appears anywhere in the list, and 0 otherwise. Therefore, unnecessary extra labels reduce the secondary score.

Additional requirements:
- If you create intermediate files, images, scripts, or notes while solving the task, create them in the current working directory.
- Internet access is disabled.
\end{DocCode}

\subsection*{\code{low_desc_code_n}}
\begin{DocCode}
Context:
The time series in the test_samples/ and train_samples/ folders were generated by simulating a ball bouncing on a surface.

Observed Signals:
col1: ball position (height)
col2: ball velocity
col3: contact impulse (N*s)
col4: time

For each simulation, either no parameter changes occur, or exactly one parameter among the allowed labels changes during the observed simulation interval.
If a parameter changes, it undergoes a single instantaneous step change at an unknown time during the observed interval.

Allowed labels:
["coefficient of restitution", "mass", "drag coefficient", "gravity acceleration", "no parameter change"]

Use "no parameter change" if there is no evidence in the data of a parameter change during the observed interval.

Task:
Create a Python script named rule.py in the current working directory.
The script must define exactly this function:
def predict(df) -> list[str]:
The input df is a pandas DataFrame containing one sample with columns col1, col2, col3, and col4.
The function will be called on samples inside test_samples/ and on additional held-out samples with the same schema and label set.
For each dataframe, predict(df) must return a ranked list of labels. The first label is the final top-1 prediction and should be the single label most likely to be correct. Additional labels are optional lower-confidence alternatives and should only be included when the evidence is genuinely ambiguous.

To help with this task, you can use the labeled train_samples/ directory while developing rule.py. The corresponding labels are available in train_labels.json.

Requirements for predict(df):
- Return a Python list of strings.
- Every returned label must exactly match one of the allowed labels.
- The first returned label is the top-1 prediction.
- Do not include duplicate labels.
- Do not return all labels unless the evidence is genuinely ambiguous across all labels.
- You may inspect train_samples/ and train_labels.json while developing rule.py, but the final submitted rule.py must not read, open, import, or depend on any files at prediction time, including train_samples/, test_samples/ or train_labels.json
- The final rule.py must be able to run on a dataframe alone.

Evaluation:
The primary evaluation metric is top-1 accuracy: the first label returned by predict(df) is compared with the hidden correct label.
A secondary shortlist score may also be reported. For a returned list of length m, the sample receives score 1/m if the hidden correct label appears anywhere in the list, and 0 otherwise. Therefore, unnecessary extra labels reduce the secondary score.

Additional requirements:
- If you create intermediate files, images, scripts, or notes while solving the task, create them in the current working directory.
- Internet access is disabled.
\end{DocCode}


\subsection*{\code{none_desc_direct_n}}
\begin{DocCode}
Context:
The time series in the test_samples/ and train_samples/ folders were generated by a simulator of an unknown physical phenomenon.
The column meanings are unknown, except for the last column, which represents time.

For each simulation, either no parameter changes occur, or exactly one parameter corresponding to one of the allowed changes during the observed simulation interval.
If a parameter changes, it undergoes a single instantaneous step change at an unknown time during the observed interval.

["label_0", "label_1", "label_2", "label_3"] denote different parameter changes, while "label_4" denotes that no parameter changed.

Task:
Create a file named results.json in the current working directory.
For each file in test_samples/, return a ranked list of labels.
The first label is your final top-1 prediction and should be the single label you think is most likely correct.
You may include additional labels only when the evidence is genuinely ambiguous.
Additional labels are treated as lower-confidence alternatives.
The output must be valid JSON with exactly this structure:
{

  "<filename_1>": ["<top_1_label>"],
  "<filename_2>": ["<top_1_label>", "<optional_lower_confidence_label>"]

}

To help with this task, you can use the labeled `train_samples/` directory. The corresponding labels are available in `train_labels.json` file.

Requirements:
- Include one entry for every Parquet file in test_samples/.
- Every returned label must exactly match one of the allowed labels.
- The order of labels matters: the first label is the top-1 prediction.
- Do not include duplicate labels for a sample.
- Do not return all labels unless the evidence is genuinely ambiguous across all labels.

Evaluation:
The primary evaluation metric is top-1 accuracy: the first label in each returned list is compared with the hidden correct label.
A secondary shortlist score may also be reported. For a returned list of length m, the sample receives score 1/m if the hidden correct label appears anywhere in the list, and 0 otherwise. Therefore, unnecessary extra labels reduce the secondary score.

Additional requirements:
- If you create intermediate files, images, scripts, or notes while solving the task, create them in the current working directory.
- Internet access is disabled.
\end{DocCode}

\subsection*{\code{none_desc_code_n}}
\begin{DocCode}
Context:
The time series in the test_samples/ and train_samples/ folders were generated by a simulator of an unknown physical phenomenon.
The column meanings are unknown, except for the last column, which represents time.

For each simulation, either no parameter changes occur, or exactly one parameter corresponding to one of the allowed changes during the observed simulation interval.
If a parameter changes, it undergoes a single instantaneous step change at an unknown time during the observed interval.

["label_0", "label_1", "label_2", "label_3"]  denote different parameter changes, while "label_4" denotes that no parameter changed.

Task:
Create a Python script named rule.py in the current working directory.
The script must define exactly this function:
def predict(df) -> list[str]:
The input df is a pandas DataFrame containing one sample with columns col1, col2, col3, and col4.
The function will be called on samples inside test_samples/ and on additional held-out samples with the same schema and label set.
For each dataframe, predict(df) must return a ranked list of labels. The first label is the final top-1 prediction and should be the single label most likely to be correct. Additional labels are optional lower-confidence alternatives and should only be included when the evidence is genuinely ambiguous.

To help with this task, you can use the labeled train_samples/ directory while developing rule.py. The corresponding labels are available in train_labels.json.

Requirements for predict(df):
- Return a Python list of strings.
- Every returned label must exactly match one of the allowed labels.
- The first returned label is the top-1 prediction.
- Do not include duplicate labels.
- Do not return all labels unless the evidence is genuinely ambiguous across all labels.
- You may inspect train_samples/ and train_labels.json while developing rule.py, but the final submitted rule.py must not read, open, import, or depend on any files at prediction time, including train_samples/, test_samples/ or train_labels.json
- The final rule.py must be able to run on a dataframe alone.

Evaluation:
The primary evaluation metric is top-1 accuracy: the first label returned by predict(df) is compared with the hidden correct label.
A secondary shortlist score may also be reported. For a returned list of length m, the sample receives score 1/m if the hidden correct label appears anywhere in the list, and 0 otherwise. Therefore, unnecessary extra labels reduce the secondary score.

Additional requirements:
- If you create intermediate files, images, scripts, or notes while solving the task, create them in the current working directory.
- Internet access is disabled.
\end{DocCode}


\tsenvAppendixPromptExampleCombinationsEnd
